\mainchapterstyle
\chapter{بحث و نتیجه‌گیری}
\pagestyle{mainheader}

\section{مقدمه}
تاکنون شما در پایان‌نامه‌ای که مشغول نوشتن آن هستید، پاسخ چهار سؤال را دادهاید:
چرا تحقیق را انجام دادید؟ (مقدمه)
دیگران در این زمینه چه کارهایی کرده‌اند و تمایز کار شما با آنها چیست؟ (مرور ادبیات)
چگونه تحقیق را انجام دادید؟ (روش‌ها)
چه از تحقیق به دست آوردید؟ (یافته‌ها)
حال زمان آن فرا رسیده که با توجه به تمامی مطالب ذکر شده، در نهایت به سؤال آخر پاسخ دهید:
چه برداشتی از یافته‌های تحقیق کردید؟ (نتیجه‌گیری)
در واقع در این بخش، هدف پاسخ به این سؤال است که چه برداشتی از یافته‌ها کردید و این یافته‌ها چه
فایده‌ای دارند؟

نتیجه‌گیری مختصری بنویسید. ارائه‌ی داده‌ها، نتایج و یافته‌ها در فصل چهارم ارائه می‌شود. در این فصل
تفاوت، تضاد یا تطابق بین نتایج تحقیق با نتایج دیگر محققان باید ذکر شود. تفسایر و تحلیل نتایج نباید بر
اساس حدس و گمان باشد، بلکه باید برمبنای نتایج عملی استخراج شده از تحقیق و یا استناد به تحقیقات
دیگران باشد. با توجه به حجم و ماهیت تحقیق و با صلاحدید استاد راهنما، این فصل میتواند تحت عنوانی
دیگر بیاید یا به دو فصل جداگانه با عناوین مناسب، تفکیک شود. این فصل فقط باید به جمع‌بندی
دست‌آوردهای فصلهای چهارم و پنجم محدود و از ذکر موارد جدید در آن خودداری شود. در عنوان این فصل، به جای کلمه تفسیر می‌توان از واژگان بحث و تحلیل هم استفاده کرد. این فصل شاید مهم‌ترین فصل پایان‌نامه باشد.

در این فصل خلاصه‌ای از یافته‌های تحقیق جاری ارائه می‌شود. این فصل میتواند حاوی یک مقدمه،
شامل مروری اجمالی بر مراحل انجام تحقیق باشد (حدود یک صفحه). مطالب پاراگراف‌بندی شود و هر
پاراگراف به یک موضوع مستقل اختصاص یابد. فقط به ارائه‌ی یافته‌ها و دست‌آوردها بسنده شود و از تعمیم
بی‌مورد نتایج خودداری شود. تا حد امکان از ارائه‌ی جداول و نمودارها اجتناب شود. از ارائه‌ی عناوین کلی در
حوزه‌ی تحقیق و پیشنهاد تحقیقات آتی خودداری شود و کاملا در چارچوب و زمینه‌ی مربوط به تحقیق جاری
باشد. این فصل حدود 10 - 15 صفحه است.

\section{محتوا}
به ترتیب شامل موارد زیر است:

\subsection{جمع‌بندی}
خلاصه‌ای از تمام یافته‌ها و دست‌آوردهای تحقیق جاری است.

\subsection{نوآوری}
نوآوری تحقیق را بر اساس یافته‌های آن تشریح می‌کند. که دارای دو بخش اصلی است: نوآوری
تئوری، یعنی تمایز تئوریک کار با کارهای محققین قبلی و نوآوری عملی، یعنی توصیه های محقق به
صنعت برای بهبود بخشیدن به کارها بر اساس یافته های تحقیق.

\subsection{پیشنهادها}
عناوین و موضوعات پیشنهادی را برای تحقیقات آتی بیشتر در زمینه‌ی مورد بحث در آینده ارائه می‌کند.


\subsection{محدودیت‌ها}
محدودیت‌هایی که کنترل آن از عهده محقق خارج است مانند انتخاب نوع یافته‌ها و دیگر محدودیت‌هایی که کنترل آن در دست محقق است مانند موضوع و محل تحقیق و ... تشریح می‌شوند. تاثیر این
محدودیت‌ها بر یافته‌های تحقیق در این قسمت شرح داده می‌شوند.